% END RX INLINE 2

\paragraph{Exact dictionary for the equilibrium provider.}
The elliptic modulus $k$ used locally in EIQ is the $\kappa$ of GEL,
and EIQ's local $m=(u^2+v^2)/2$ is an endpoint midpoint. They are
distinctly indexed coordinates from the original packet
$k_{\rm packet}=4l+1$ and multiplicity $m_{\rm packet}$.
GEL3 gives their exact parameter map
$\alpha=a/n$, $\beta=\pi q/(2n)$, tending to $(2,\pi)$ on
the original sequences. GEL's constant $C_\Gamma=\Gamma(1/4)^2/(2\pi)$
in its source bound retains that definition; the leading coefficient
is the separately defined $\mathcal C_\Gamma$ in GEL18.

% BEGIN RX INLINE 3; SHA256 a0a0346a19d0e035a186590dbd69881cdf7a30e8fef602d47bba0932f175e731
% Independent proof. No external one-cut or equilibrium theorem is used.
% May be included as a section in the cumulative paper.
\section{Exact equilibrium for the square-root logarithmic potential}
\label{sec:independent-sqrt-log-equilibrium}

Throughout this section $\alpha>0$ and $\beta>0$.  On $(0,\infty)$ set
\begin{equation}\tag{EIQ1}
 V_{\alpha,\beta}(x)=\beta\sqrt{x}-\alpha\log x,
 \qquad
 \mathcal E_{\alpha,\beta}(\mu)
 =\int V_{\alpha,\beta}(x)\,d\mu(x)
       -\iint\log|x-y|\,d\mu(x)d\mu(y).
\end{equation}
The measures considered initially are probability measures for which
$\int V_{\alpha,\beta}\,d\mu$ is finite.  Because this potential is bounded
below and tends to $+\infty$ at both ends, this implies
$\int(\sqrt{x}+|\log x|)\,d\mu<\infty$.
The positive part of $\log|x-y|$ is then integrable, by
$\log^+|x-y|\leq\log(1+x)+\log(1+y)$.
Thus the energy is well defined with value in $\mathbb R\cup\{+\infty\}$.
All comparisons below also apply to the usual extended energy defined by
integrating the kernel
$\tfrac12V_{\alpha,\beta}(x)+\tfrac12V_{\alpha,\beta}(y)-\log|x-y|$:
the kernel is bounded below, and finite energy for this kernel implies the
same moment conditions.  To see the latter assertion without cancellation,
put $g(x)=\tfrac12V_{\alpha,\beta}(x)-\log(1+x)$.
The kernel is at least $g(x)+g(y)$; $g$ is bounded below and controls positive
constant multiples of both $\sqrt{x}$ near infinity and $|\log x|$ near zero.

\subsection{Existence and uniqueness of the two endpoints}
For $0<k<1$ define the complete elliptic integrals by their actual real
integrals
\begin{equation}\tag{EIQ2}
 K(k)=\int_0^{\pi/2}\frac{d\theta}
            {\sqrt{1-k^2\sin^2\theta}},\qquad
 E(k)=\int_0^{\pi/2}\sqrt{1-k^2\sin^2\theta}\,d\theta.
\end{equation}
There is exactly one $k\in(0,1)$ satisfying
\begin{equation}\tag{EIQ3}
 \frac{\sqrt{1-k^2}K(k)}{E(k)}=\frac{\alpha}{\alpha+2}.
\end{equation}
For this $k$ put
\begin{equation}\tag{EIQ4}
 r=\sqrt{1-k^2},\qquad
 v=\frac{(\alpha+2)\pi}{\beta E(k)},\qquad u=rv,
 \qquad a=u^2,\quad b=v^2.
\end{equation}
In particular $0<a<b<\infty$, and the endpoint equations, with all
constants retained, are
\begin{equation}\tag{EIQ5}
 uK(k)=\frac{\alpha\pi}{\beta},\qquad
 vE(k)=\frac{(\alpha+2)\pi}{\beta},\qquad
 k^2=1-\frac{u^2}{v^2}.
\end{equation}

Here is a proof of the assertion of uniqueness. Differentiation of the
integrals in (EIQ2), and integration of the derivative of
$\sin\theta\cos\theta/\sqrt{1-k^2\sin^2\theta}$ between its two zero
endpoint values, give
\[
 E'(k)=\frac{E(k)-K(k)}{k},\qquad
 K'(k)=\frac{E(k)}{k(1-k^2)}-\frac{K(k)}{k}.
\]
For completeness the second identity is equivalent, after differentiation
under the integral sign, to
\[
 \int_0^{\pi/2}
 \left\{\frac{k^2\sin^2\theta}{(1-k^2\sin^2\theta)^{3/2}}
 -\frac{1-k^2\sin^2\theta}{(1-k^2)\sqrt{1-k^2\sin^2\theta}}
 +\frac1{\sqrt{1-k^2\sin^2\theta}}\right\}d\theta=0;
\]
the expression in braces is
$-k^2(1-k^2)^{-1}\,d(\sin\theta\cos\theta/
\sqrt{1-k^2\sin^2\theta})/d\theta$.
With $t=E(k)/K(k)$, the integral definitions show strictly
$r^2<t<1$.  Consequently, if $f(k)=rK(k)/E(k)$, then
\begin{equation}\tag{EIQ6}
 \frac{kf'(k)}{f(k)}
 =\frac{t}{r^2}+\frac1t-1-\frac1{r^2}
 =\frac{(t-1)(t-r^2)}{r^2t}<0.
\end{equation}
At $k=0$ the continuous limiting value is $f(0)=1$.
As $k\uparrow1$, the integrand of $rK(k)$ tends to zero at every
$\theta<\pi/2$ and is bounded by $1$, while $E(k)\to1$ by dominated
convergence.  Hence $f(k)\to0$.  This proves (EIQ3), including existence.
It also proves smooth dependence of $k,u,v,a,b$ on $(\alpha,\beta)$ in
$(0,\infty)^2$, by the nonzero derivative in (EIQ6).

Let $m=(a+b)/2$, $R_s=\sqrt{(a+s)(b+s)}$ for $s\geq0$, and
$R_0=uv$.  Substitution
$t=a\cos^2\theta+b\sin^2\theta$ gives, directly from (EIQ5),
\begin{equation}\tag{EIQ7}
 \int_a^b\frac{V'_{\alpha,\beta}(t)}
                  {\sqrt{(b-t)(t-a)}}\,dt=0,
 \qquad
 \int_a^b\frac{tV'_{\alpha,\beta}(t)}
                  {\sqrt{(b-t)(t-a)}}\,dt=2\pi.
\end{equation}
Indeed the three integrals with numerators $t^{-1/2},t^{-1},t^{1/2}$
are respectively $2K(k)/v$, $\pi/(uv)$, and $2vE(k)$, and
$V'(t)=\beta/(2\sqrt t)-\alpha/t$.

\subsection{A positive density and its exact Cauchy transform}
The elementary Stieltjes integral, obtained by $s=xz^2$, is
\begin{equation}\tag{EIQ8}
 \int_0^\infty\frac{ds}{\sqrt s(x+s)}=\frac\pi{\sqrt x}
 \quad(x>0).
\end{equation}
Using this identity in (EIQ7) and applying Tonelli to its positive
constituents gives
\begin{equation}\tag{EIQ9}
 \frac\alpha{R_0}=
 \frac\beta{2\pi}\int_0^\infty\frac{ds}{\sqrt s R_s},
 \qquad
 \int_0^\infty\frac{1-s/R_s}{\sqrt s}\,ds=2vE(k).
\end{equation}
Both integrals converge; their integrands are $O(s^{-1/2})$ at zero
and $O(s^{-3/2})$ at infinity.  The second formula also follows by
averaging $\sqrt t=\pi^{-1}\int_0^\infty
s^{-1/2}t/(t+s)\,ds$ against the arcsine probability measure on $[a,b]$.

Define for every $x>0$
\begin{align}\tag{EIQ10}
 h(x)
 &=\frac\alpha{xR_0}
   -\frac\beta{2\pi}
       \int_0^\infty\frac{ds}{\sqrt s(x+s)R_s} \\
 &=\frac\beta{2\pi x}
       \int_0^\infty\frac{\sqrt s}{(x+s)R_s}\,ds>0.\notag
\end{align}
The second equality is exactly (EIQ9), with no change of the potential.
All integrals in (EIQ10) converge absolutely and locally uniformly in
$x>0$.  The candidate density on the original positive half-line is
\begin{equation}\tag{EIQ11}
 \rho_{\alpha,\beta}(x)=
 \begin{cases}
 \displaystyle
 \frac{\sqrt{(b-x)(x-a)}}{2\pi}\,h(x),&a<x<b,\\
 0,&x\notin(a,b).
 \end{cases}
\end{equation}
It is strictly positive on $(a,b)$ and vanishes continuously with the
square-root factor at each endpoint.  Positivity here follows from the
last integral in (EIQ10), not from a one-interval ansatz.

The following elementary transform computes its mass and resolvent.
Choose $R(z)=\sqrt{(z-a)(z-b)}$ on $\mathbb C\setminus[a,b]$ with
$R(z)/z\to1$ at infinity. Thus $R(x)>0$ for $x>b$,
$R(x)<0$ for $x<a$, and $R(x+i0)=i\sqrt{(b-x)(x-a)}$ on $(a,b)$.
Direct substitution $t=m+(b-a)\cos\theta/2$, followed by
$w=\tan(\theta/2)$, proves
\[
 \frac1\pi\int_a^b\frac{dt}{(z-t)\sqrt{(b-t)(t-a)}}=\frac1{R(z)}.
\]
For real $z>b$, the substitution reduces the integral to
$\pi^{-1}\int_0^\infty2\,dw/[(z-b)+(z-a)w^2]$;
the identity elsewhere follows by analytic continuation.
Polynomial division then gives
\begin{equation}\tag{EIQ12}
 \frac1\pi\int_a^b\frac{\sqrt{(b-t)(t-a)}}{z-t}\,dt
     =z-m-R(z),\qquad
 \frac1\pi\int_a^b\frac{\sqrt{(b-t)(t-a)}}{t+s}\,dt
     =m+s-R_s.
\end{equation}
For example the first identity follows from
$(b-t)(t-a)=-(t-z)^2+2(m-z)(t-z)-(z-a)(z-b)$,
the preceding arcsine transform, and the arcsine first moment $m$.
Partial fractions in these two identities yield
\begin{equation}\tag{EIQ13}
 H_s(z):=\frac1\pi\int_a^b
       \frac{\sqrt{(b-t)(t-a)}}{(t+s)(z-t)}\,dt
       =1-\frac{R_s+R(z)}{z+s}.
\end{equation}
The formula at $z=-s$ is interpreted by its removable limit.

The mass of (EIQ11), computed using (EIQ12), is
\begin{align}\tag{EIQ14}
 \int_a^b\rho_{\alpha,\beta}(t)\,dt
 &=\frac\alpha{2R_0}(m-R_0)
 -\frac\beta{4\pi}\int_0^\infty
       \frac{(m+s)/R_s-1}{\sqrt s}\,ds\\
 &=-\frac\alpha2+
   \frac\beta{4\pi}\int_0^\infty\frac{1-s/R_s}{\sqrt s}\,ds
 =-\frac\alpha2+\frac{\beta vE(k)}{2\pi}=1.\notag
\end{align}
In particular this is a probability density.

Define its Cauchy transform with the explicit sign convention
$G(z)=\int_a^b\rho_{\alpha,\beta}(t)/(z-t)\,dt$.
For $z$ outside both $[a,b]$ and $(-\infty,0]$, formulas (EIQ10) and
(EIQ13) give
\begin{align}\tag{EIQ15}
 G(z)
 &=\frac\alpha{2R_0}H_0(z)
       -\frac\beta{4\pi}
          \int_0^\infty\frac{H_s(z)}{\sqrt sR_s}\,ds\\
 &=-\frac\alpha{2z}+\frac\beta{4\sqrt z}
       -\frac{R(z)h(z)}2
 =\frac{V'_{\alpha,\beta}(z)-R(z)h(z)}2.\notag
\end{align}
Here the positive square root is continued with its negative-real-axis
cut, and $h(z)$ is the first integral formula of (EIQ10).  Every integral
exchange is absolutely convergent, uniformly on compact subsets of this
domain: near zero the auxiliary integrand is $O(s^{-1/2})$, and near
infinity the original integral defining $H_s$ is $O(s^{-1})$.
The constant terms on the second line cancel by (EIQ9), and (EIQ8)
supplies the term $\beta/(4\sqrt z)$.

\subsection{The Euler equality, its strict exterior inequality, and uniqueness}
Set
\begin{equation}\tag{EIQ16}
 U(x)=V_{\alpha,\beta}(x)
             -2\int_a^b\log|x-t|\rho_{\alpha,\beta}(t)\,dt
 \quad(x>0).
\end{equation}
This function is continuous on $(0,\infty)$.  The integral is continuous
because the density is bounded, compactly supported, and the integral
of $|\log|x-t||$ over an interval of length $\varepsilon$ tends to zero
uniformly in its centre.  In the interior of $[a,b]$, differentiation of
the logarithmic potential is the principal-value integral
$\operatorname{PV}\int\rho(t)/(x-t)\,dt$.
This follows by subtracting $\rho(x)$ on a symmetric interval about
$x$; the remaining numerator is $O(|t-x|)$ because the density is
continuously differentiable locally in the interior.  The real part of
(EIQ15) on the upper edge of the interval is precisely this principal
value.  Indeed subtracting $\rho(x)$ in
$\int\rho(t)/(x+i\varepsilon-t)\,dt$ gives the same limit, while the
constant term has real part equal to its elementary principal value.
Since $R(x+i0)$ is purely imaginary and $h(x)$ is real, (EIQ15) proves
$U'(x)=0$ on $(a,b)$.  There is therefore a real constant $\ell$ such that
\begin{equation}\tag{EIQ17}
 U(x)=\ell\quad(a\leq x\leq b).
\end{equation}
For positive $x$ outside this interval, (EIQ15) instead gives the exact
real derivative
\begin{equation}\tag{EIQ18}
 U'(x)=R(x)h(x)
 \begin{cases}<0,&0<x<a,\\>0,&x>b.\end{cases}
\end{equation}
Combining (EIQ17), continuity, and these signs proves
\begin{equation}\tag{EIQ19}
 U(x)>\ell\quad\hbox{for }x\in(0,a)\cup(b,\infty).
\end{equation}
This verifies the entire exterior inequality on the original domain,
including the interval next to zero.

Here is a complete positivity argument for the energy of the difference
of two measures.  If $\nu$ is a finite signed real measure of total mass
zero for which $|\log|x-y||$ is integrable against
$|\nu|\otimes|\nu|$, the scalar identity
\begin{equation}\tag{EIQ20}
 -\log d=\frac12\int_0^\infty
           \frac{e^{-td^2}-e^{-t}}t\,dt\quad(d>0)
\end{equation}
follows by differentiating in $d$ and evaluating at $d=1$.
One half of the integral over any finite subinterval of $(0,\infty)$
has the same sign as $-\log d$ and absolute value at most $|\log d|$.
Consequently dominated convergence and the zero total mass give
\begin{equation}\tag{EIQ21}
 -\iint\log|x-y|\,d\nu(x)d\nu(y)
 =\frac12\int_0^\infty\frac1t
       \iint e^{-t(x-y)^2}\,d\nu(x)d\nu(y)\,dt\geq0.
\end{equation}
The nonnegative sign is proved, rather than assumed, by the Gaussian
Fourier identity
\begin{equation}\tag{EIQ22}
 \iint e^{-t(x-y)^2}\,d\nu(x)d\nu(y)
 =\frac1{\sqrt{4\pi t}}\int_{\mathbb R}
          e^{-\xi^2/(4t)}|\widehat\nu(\xi)|^2\,d\xi,
 \qquad \widehat\nu(\xi)=\int e^{i\xi x}\,d\nu(x).
\end{equation}
The scalar Gaussian Fourier identity follows by completing the square
in the Gaussian integral (or by differentiating its Fourier transform
and its value at zero); its integration against the finite measures is
justified by absolute integrability.  Equality in (EIQ21) forces
$\widehat\nu=0$: the right side of (EIQ22) is strictly positive for every
$t>0$ if this continuous Fourier transform has a nonzero value.
Finally $\widehat\nu=0$ implies $\nu=0$ without an extra uniqueness
assumption. Convolve $\nu$ with a centred Gaussian approximate identity.
Its value at every point is zero by that same Fourier integral.
For every compactly supported continuous test function, its convolution
with these Gaussians converges uniformly to the test function; integration
against $\nu$ then proves that $\nu$ annihilates every such function.

Let $\mu$ be any probability measure of finite energy in (EIQ1), and
write $\mu_*(dx)=\rho_{\alpha,\beta}(x)\,dx$ and $\nu=\mu-\mu_*$.
The integrability needed for (EIQ21) holds.  The $\mu\otimes\mu$ term
has it by finite energy and the already established logarithmic moment;
the $\mu_*\otimes\mu_*$ term has it by bounded compact support; and the
cross term has it because the negative logarithmic singularity integrated
against the bounded density is uniformly bounded, while the positive
part is controlled by $\log(1+x)+\log(1+y)$.
Expanding the energy therefore yields exactly
\begin{equation}\tag{EIQ23}
 \mathcal E_{\alpha,\beta}(\mu)-
       \mathcal E_{\alpha,\beta}(\mu_*)
 =\int(U(x)-\ell)\,d\mu(x)
       -\iint\log|x-y|\,d\nu(x)d\nu(y)\geq0.
\end{equation}
The first term is nonnegative by (EIQ17)--(EIQ19), and the second is
strictly positive unless $\mu=\mu_*$. Infinite energy competitors cannot
improve this finite value. Thus (EIQ11) is the unique global minimizing
probability measure, with exactly the support $[a,b]$.

\subsection{An explicit endpoint integral for the logarithmic moment}
Retain $m,R_s,u,v$ as defined above, and define
\begin{equation}\tag{EIQ24}
 C=\frac{(u+v)^2}4,\qquad
 \eta=\frac{v-u}{v+u},\qquad
 \zeta_s=\frac{b-a}{2(m+s+R_s)}\quad(s\geq0).
\end{equation}
Thus $0<\zeta_s\leq\eta<1$ and $\zeta_0=\eta$.
For every $s\geq0$ put
\begin{equation}\tag{EIQ25}
 L_s=(m+s-R_s)\log C-(m-R_0)
                  -2R_s\log(1-\eta\zeta_s).
\end{equation}
Then the logarithmic moment of the equilibrium measure is exactly
\begin{equation}\tag{EIQ26}
 \mathcal L(\alpha,\beta):=\int_a^b\log x\rho_{\alpha,\beta}(x)\,dx
 =\frac\beta{4\pi}\int_0^\infty
                    \frac{L_0-L_s}{\sqrt s R_s}\,ds.
\end{equation}
This is a convergent one-dimensional endpoint integral of explicitly
given real elementary functions; its only endpoints are those specified
uniquely in (EIQ3)--(EIQ5).

Here is a direct derivation, including the constants in (EIQ25).
Set $h_0=(b-a)/2$ and substitute $t=m+h_0\cos\theta$ in arcsine
integrals.  The identities
\[
 t=C(1+2\eta\cos\theta+\eta^2),\qquad
 t+s=\frac{m+s+R_s}2
                (1+2\zeta_s\cos\theta+\zeta_s^2)
\]
give the absolutely convergent Fourier series
\[
 \log t=\log C+
       2\sum_{j\geq1}\frac{(-1)^{j+1}\eta^j}j\cos(j\theta),
 \qquad
 \frac1{t+s}=\frac1{R_s}
       \left(1+2\sum_{j\geq1}(-\zeta_s)^j\cos(j\theta)\right).
\]
Orthogonality on $[0,\pi]$ gives
\begin{align}\tag{EIQ27}
 \frac1\pi\int_a^b\frac{\log t}{\sqrt{(b-t)(t-a)}}\,dt
 &=\log C,\\
 \frac1\pi\int_a^b\frac{t\log t}{\sqrt{(b-t)(t-a)}}\,dt
 &=m\log C+(m-R_0),\notag\\
 \frac1\pi\int_a^b\frac{\log t}{(t+s)\sqrt{(b-t)(t-a)}}\,dt
 &=\frac{\log C+2\log(1-\eta\zeta_s)}{R_s}.\notag
\end{align}
The second identity uses $h_0\eta=m-R_0$.
Polynomial division
\[
 \frac{(b-t)(t-a)}{t+s}
       =-t+2m+s-\frac{R_s^2}{t+s}
\]
now proves, exactly,
\begin{equation}\tag{EIQ28}
 L_s=\frac1\pi\int_a^b
               \frac{\log t\sqrt{(b-t)(t-a)}}{t+s}\,dt.
\end{equation}
Using the positive integral expression for $h$ in (EIQ10), and the
identity $1/[t(t+s)]=(1/t-1/(t+s))/s$, yields (EIQ26).
The exchanges are absolutely convergent: $\log t$ is bounded on
$[a,b]$, $L_0-L_s=O(s)$ at zero by (EIQ28), and $L_s=O(s^{-1})$
at infinity.  Thus the integrand in (EIQ26) is $O(s^{1/2})$ at zero
and $O(s^{-3/2})$ at infinity.

\subsection{Parameter continuity and the variational derivative}
The density and its logarithmic moment depend smoothly on
$(\alpha,\beta)\in(0,\infty)^2$ in the following precise sense.
On a compact parameter neighbourhood, (EIQ3)--(EIQ6) put $a,b$ in a
fixed compact subset of $(0,\infty)$, with $b-a$ bounded below by a
positive number.  The integral for $h$ and each of its parameter and
$x$ derivatives converge uniformly for $x$ in this fixed compact set:
their absolute integrands are bounded by constants times $s^{1/2}$
near zero and $s^{-3/2}$ near infinity. Parameter derivatives of the
endpoints are bounded on a smaller compact neighbourhood.
Under the fixed-interval coordinate
$x=a+(b-a)y$, $0\leq y\leq1$, the measure is exactly
\begin{equation}\tag{EIQ29}
 \rho_{\alpha,\beta}(x)\,dx
 =\frac{(b-a)^2}{2\pi}\sqrt{y(1-y)}
                  h\bigl(a+(b-a)y\bigr)\,dy.
\end{equation}
All parameter derivatives of its coefficient are bounded by a constant
times $\sqrt{y(1-y)}$. The same is true after multiplication by
$\log(a+(b-a)y)$ or by $\sqrt{a+(b-a)y}$.
Differentiation under this integral proves smoothness, in particular
continuity, of both moments on the entire positive parameter quadrant.

Define $E_{\min}(\alpha,\beta)=\mathcal E_{\alpha,\beta}(\mu_{\alpha,\beta})$,
where the measure is the unique minimizer proved above.
Testing the minimizer for a parameter pair against the neighbouring
potential, and then reversing the two tests, gives for $h>0$
\[
 -\mathcal L(\alpha+h,\beta)
 \leq\frac{E_{\min}(\alpha+h,\beta)-E_{\min}(\alpha,\beta)}h
 \leq-\mathcal L(\alpha,\beta).
\]
The corresponding inequalities for $h<0$ and the proved continuity
therefore give
\begin{equation}\tag{EIQ30}
 \partial_\alpha E_{\min}(\alpha,\beta)=-\mathcal L(\alpha,\beta),\qquad
 \partial_\beta E_{\min}(\alpha,\beta)
        =\int\sqrt{x}\,d\mu_{\alpha,\beta}(x).
\end{equation}
No unproved differentiability of a minimizing selection is being used.
For the opposite convention
$F(\alpha,\beta)=\sup_\mu\{\iint\log|x-y|\,d\mu(x)d\mu(y)
-\int V_{\alpha,\beta}\,d\mu\}=-E_{\min}(\alpha,\beta)$,
the exact derivative is therefore
$\partial_\alpha F(\alpha,\beta)=\mathcal L(\alpha,\beta)$.

The second moment in (EIQ30) has an exact value.  Push the minimizing
measure forward by the explicit dilation $x\mapsto c^2x$, $c>0$.
Writing $M=\int\sqrt{x}\,d\mu_{\alpha,\beta}(x)$, its energy minus
the energy at $c=1$ is precisely
$\beta(c-1)M-2(\alpha+1)\log c$.
Its derivative at the minimizing value $c=1$ is zero. Hence
\begin{equation}\tag{EIQ31}
 \int\sqrt{x}\,d\mu_{\alpha,\beta}(x)
       =\frac{2(\alpha+1)}\beta.
\end{equation}
The same explicit dilation, or (EIQ3)--(EIQ5), proves that the map
$x\mapsto(\pi/\beta)^2x$ pushes $\mu_{\alpha,\pi}$ to
$\mu_{\alpha,\beta}$, including the density and its two endpoints.
Consequently
\begin{equation}\tag{EIQ32}
 \mathcal L(\alpha,\beta)
       =\mathcal L(\alpha,\pi)+2\log(\pi/\beta).
\end{equation}

For the particular original limiting potential
$V(x)=\pi\sqrt{x}-2\log x$, the substitutions are exactly
$\alpha=2$, $\beta=\pi$. Equations (EIQ5), (EIQ10), and (EIQ26) become
\begin{equation}\tag{EIQ33}
 uK(k)=2,\qquad vE(k)=4,\qquad
 h(x)=\frac1{2x}\int_0^\infty\frac{\sqrt s}{(x+s)R_s}\,ds,
 \qquad
 \mathcal L(2,\pi)=\frac14\int_0^\infty
                         \frac{L_0-L_s}{\sqrt sR_s}\,ds.
\end{equation}
In particular the preceding proof includes positivity, unit mass,
the full Euler inequality, the unique global minimizer, and the
parameter continuity needed at these actual parameters.

% END RX INLINE 3


% BEGIN RX INLINE 4; SHA256 2014e7fefdcd963b70b00fb89d597e24c0826582a2541dd1ff44d63051979fc1
% Original-source large-degree calculation, 14 September 2026.
% Input fragment. Equilibrium provider:
% independent/EXACT_SQRT_LOG_EQUILIBRIUM.tex (EIQ1--33).
\section{The exact Gamma ensemble and the leading growing return}
\label{sec:GEL}

\paragraph{GEL1. Source, exponents, and the literal coordinate map.}
The source throughout is
\[
 d\sigma(y)=\frac{|\Gamma(1/4+iy/2)|^2}{2\pi}\,dy,
 \qquad \sigma(\mathbb R)=\sqrt{2\pi},\qquad
 d\nu(t)=\frac{|\Gamma(1/4+i\sqrt t/2)|^2}{2\pi\sqrt t}\,dt.
 \tag{GEL1}
\]
Here \(\nu\) is the full pushforward under \(y\mapsto y^2\); both
half-axes are present in its density. For integers \(a,n\geq0\), retain
\[
 H_n^{(a)}=\det\left[\int_0^\infty t^{a+i+j}\,d\nu(t)\right]_{i,j<n}.
\]
The definition extends to real \(a>-1/2\); all its integrals converge.
For \(n\geq1\), expansion of the two Vandermonde determinants followed
by integration term by term gives the exact identity
\[
 H_n^{(a)}=\frac1{n!}\int_{(0,\infty)^n}
       \prod_{i<j}(t_i-t_j)^2\prod_{i=1}^n t_i^a\,d\nu(t_i).
 \tag{GEL2}
\]
Indeed expansion gives \(n!\) copies of the Gram determinant after one
of the two permutation indices is fixed. Absolute integrability follows
by expanding the squared polynomial and using the existing moments.

Fix a positive real coordinate scale \(q\), independently of \(a,n\),
and make the literal map \(t_i=q^2x_i\). Introduce the reference
probability measure
\[
 d\omega(x)=\tfrac12x^{-1/2}e^{-\sqrt x}\,dx,
 \qquad B_q(x)=2q e^{\sqrt x}\sigma(q\sqrt x).
\]
Its mass is one by \(x=u^2\). This reference measure does not replace
or rescale the source: the exact Radon--Nikodym identity is
\(d\nu(q^2x)=B_q(x)d\omega(x)\), including its whole mass. Consequently
\[
 \boxed{H_n^{(a)}=q^{2an+2n(n-1)}J_{n,q}(a),\quad
 J_{n,q}(a)=\frac1{n!}\int\prod_{i<j}(x_i-x_j)^2
                         \prod_i x_i^aB_q(x_i)\,d\omega(x_i).}
 \tag{GEL3}
\]
The powers in this identity come separately from the \(n\) factors
\(t_i^a\) and the \(n(n-1)/2\) squared differences. No determinant,
source mass, or Jacobian has been omitted.

\paragraph{GEL2. An elementary global Gamma bound with the exact exponential.}
Put \(\lambda=1/4\), \(b\geq0\), and
\(C_\Gamma=\Gamma(1/4)^2/(2\pi)\). Euler's finite limit gives
\[
 \frac{|\Gamma(\lambda+ib)|^2}{\Gamma(\lambda)^2}
 =\prod_{j=0}^\infty\left(1+\frac{b^2}{(j+\lambda)^2}\right)^{-1}.
\]
For clarity, the finite limit follows by applying the beta integral to
\(\int_0^1t^{z-1}(1-t)^Ndt=N!/[z(z+1)\cdots(z+N)]\), replacing
\(t\) by \(u/N\), and using dominated convergence for
\(\operatorname{Re}z>0\). Taking squared moduli and dividing by the
same limit at \(z=\lambda\) proves the displayed product. Its logarithm
converges absolutely because its summand is \(O((j+1)^{-2})\).

The decreasing function \(v\mapsto\log(1+b^2/v^2)\) gives, with
\(S=\sum_{j\geq0}\log(1+b^2/(j+\lambda)^2)\),
\[
 I\leq S\leq I+\log(1+b^2/\lambda^2),\qquad
 I=\int_\lambda^\infty\log(1+b^2/v^2)\,dv
  =\pi b-\lambda\log(1+b^2/\lambda^2)
                      -2b\arctan(\lambda/b).
\]
The integral is evaluated by differentiating
\(v\log(1+b^2/v^2)+2b\arctan(v/b)\); at \(b=0\) it is zero
by continuity. Since \(0\leq2b\arctan(\lambda/b)\leq2\lambda\),
substitution of \(b=|y|/2\) proves the global, finite inequalities
\[
 \boxed{C_\Gamma(1+4y^2)^{-3/4}e^{-\pi|y|/2}
 \leq\sigma(y)\leq
 C_\Gamma e^{1/2}(1+4y^2)^{1/4}e^{-\pi|y|/2}.}
 \tag{GEL4}
\]
Both inequalities include \(y=0\). In particular, for
\(C_0=\Gamma(1/4)^2/\pi\),
\[
 B_q(x)\leq C_0q e^{1/2}(1+2q)^{1/2}
            e^{-(\pi q/2-3/2)\sqrt x}\quad(x>0).
 \tag{GEL5}
\]
Here \((1+4q^2x)^{1/4}\leq(1+2q\sqrt x)^{1/2}
\leq(1+2q)^{1/2}(1+\sqrt x)^{1/2}
\leq(1+2q)^{1/2}e^{\sqrt x/2}\). On each fixed compact interval
\([A,B]\subset(0,\infty)\), the lower bound is
\[
 B_q(x)\geq C_0q(1+4q^2B)^{-3/4}e^{-\pi q\sqrt x/2}.
 \tag{GEL6}
\]
Thus the exact exponential, and the size and position of the remaining
source factor, have been proved directly from the original Gamma function.

\paragraph{GEL3. The energy object and its original parameter domain.}
For \(\alpha,\beta>0\) put
\[
 V_{\alpha,\beta}(x)=\beta\sqrt x-\alpha\log x,
 \quad K_{\alpha,\beta}(x,y)=\log|x-y|
           -\tfrac12V_{\alpha,\beta}(x)-\tfrac12V_{\alpha,\beta}(y),
\]
\[
 \mathcal F(\alpha,\beta)=\sup_{\mu\in\mathcal P((0,\infty))}
                   \iint K_{\alpha,\beta}(x,y)\,d\mu(x)d\mu(y).
 \tag{GEL7}
\]
The kernel is bounded above, so its integral has an unambiguous value
in \(\mathbb R\cup\{-\infty\}\) for every probability in this domain.
Whenever its value is finite it equals
\(\iint\log|x-y|\,d\mu^2-\int V_{\alpha,\beta}\,d\mu\).
Indeed \(K\leq r(x)+r(y)\), with
\(r=\log(1+x)-V_{\alpha,\beta}/2\) bounded above and tending to
\(-\infty\) at both endpoints; finite kernel integral therefore forces
the \(\sqrt x\) and \(|\log x|\) moments to be finite. The positive
logarithmic part is then integrable, and finite kernel integral forces
the negative part to be integrable as well. This defines the extended
energy without forming an indeterminate difference of infinite quantities.
The equilibrium calculation supplied with this section proves existence,
uniqueness, and the full exterior variational inequality for its density
\(\rho_{\alpha,\beta}\), with support \([u^2,v^2]\). Its exact endpoints
and positive density are
\[
 uK(\kappa)=\frac{\alpha\pi}{\beta},\quad
 vE(\kappa)=\frac{(\alpha+2)\pi}{\beta},\quad
 \kappa^2=1-u^2/v^2,
\]
\[
 \rho_{\alpha,\beta}(x)=
 \frac{\beta\sqrt{(v^2-x)(x-u^2)}}{4\pi^2x}
 \int_0^\infty
 \frac{\sqrt s\,ds}{(x+s)\sqrt{(u^2+s)(v^2+s)}}
       \quad(u^2<x<v^2).
 \tag{GEL8}
\]
The endpoint proof shows continuous parameter dependence, finite entropy
relative to \(\omega\), finite logarithmic energy, and
\[
 \partial_\alpha\mathcal F(\alpha,\beta)
       =\int\log x\,\rho_{\alpha,\beta}(x)\,dx.
 \tag{GEL9}
\]
In particular \(\mathcal F\) and this derivative are continuous on the
positive parameter quadrant. The following ensemble argument proves the
connection to GEL3; it does not presume a random-matrix asymptotic theorem.

\paragraph{GEL4. Complete free-energy limit from a compactified kernel.}
Let \(n\to\infty\), \(q/n\to\theta>0\), and let \(\alpha>0\) be fixed.
Then the actual integral in GEL3 satisfies
\[
 \boxed{\frac1{n^2}\log J_{n,q}(n\alpha)
       \longrightarrow\mathcal F(\alpha,\pi\theta/2).}
 \tag{GEL10}
\]
Here is a proof of both bounds. For the upper bound GEL5 removes the
factor \([C_0q e^{1/2}(1+2q)^{1/2}]^n/n!\), whose logarithm divided by
\(n^2\) tends to zero. The remaining exponent is
\[
 2\sum_{i<j}\log|x_i-x_j|-n\sum_iV_{\alpha,\beta_n}(x_i),
 \qquad \beta_n=\frac{\pi q}{2n}-\frac3{2n}.
\]
For \(n>1\) it is exactly
\(\sum_{i\ne j}K_{Q_n}(x_i,x_j)\), where
\[
 Q_n=\frac n{n-1}V_{\alpha,\beta_n},\qquad
 K_Q(x,y)=\log|x-y|-\tfrac12Q(x)-\tfrac12Q(y).
\]
Write \(\beta=\pi\theta/2\). Choose \(\varepsilon>0\) so that
\(\alpha-\varepsilon>0\) and \(\beta-3\varepsilon>0\).
For all sufficiently large \(n\), the two coefficients of \(Q_n\)
differ from \((\alpha,\beta)\) by at most \(\varepsilon\). It follows
pointwise that
\(Q_n\geq Q_-:=V_{\alpha-\varepsilon,\beta-3\varepsilon}\).
Indeed its difference has the form
\(b\sqrt x-a\log x\), where \(b\geq2\varepsilon\) and
\(0\leq a\leq2\varepsilon\); it is nonnegative for \(x<1\), and for
\(x\geq1\) use \(\log x\leq\sqrt x\).

Compactify \((0,\infty)\) by adjoining its two endpoints.
The function \(K_{Q_-}\), assigned value \(-\infty\) on the diagonal
and at either added endpoint, has continuous truncations
\(K_M=\max\{K_{Q_-},-M\}\). To verify the endpoint assertion uniformly
in the other variable, use
\[
 K_{Q_-}(x,y)\leq r(x)+r(y),\qquad
 r(x)=\log(1+x)-Q_-(x)/2.
\]
The function \(r\) is bounded above and tends to \(-\infty\) at both
endpoints. Continuity at a diagonal point in the interior follows from
\(\log|x-y|\to-\infty\) and local boundedness of \(Q_-\).
For the empirical probability \(L_n=n^{-1}\sum_i\delta_{x_i}\),
\[
 \frac1{n^2}\sum_{i\ne j}K_{Q_-}(x_i,x_j)
 \leq\iint K_M\,dL_n^2+\frac Mn
 \leq s_M+\frac Mn,
 \quad s_M=\max_{\mu}\iint K_M\,d\mu^2.
\]
Since \(\omega\) has mass one, integration preserves this upper bound.
Moreover \(s_M\downarrow\mathcal F(\alpha-\varepsilon,
\beta-3\varepsilon)\). For completeness, choose maximizing probabilities
\(\mu_M\). Compactness gives a convergent subsequence. For every fixed
\(L\), the inequality \(K_M\leq K_L\) when \(M\geq L\), followed by
continuity of the integral of \(K_L\), bounds the subsequential limit
of \(s_M\) by \(\iint K_Ld\mu^2\). Decrease \(L\)'s truncation to the
untruncated kernel. Monotone convergence after subtracting a common upper
bound gives \(\lim s_M\leq\iint K_{Q_-}d\mu^2\leq\mathcal F\).
The reverse inequality follows by testing every \(K_M\) on an arbitrary
probability of finite energy. Probabilities charging an added endpoint
have energy \(-\infty\), so they do not alter the supremum.
First take \(n\to\infty\), then \(M\to\infty\), and finally
\(\varepsilon\downarrow0\). Continuity proved in GEL9 yields the
upper bound in GEL10.

For the lower bound choose exactly
\(d\mu(x)=\rho_{\alpha,\beta}(x)dx\), supported on \([A,B]\), and use
GEL6. Write \(p=d\mu/d\omega\). Restricting the integral to this support
and changing measure to \(\mu^n\), Jensen's inequality gives
\begin{align*}
 \log J_{n,q}(n\alpha)\geq{}&-
 \log n!+n\log\{C_0q(1+4q^2B)^{-3/4}\}\\
 &+n(n-1)\iint\log|x-y|\,d\mu(x)d\mu(y)\\
 &-n^2\int V_{\alpha,\pi q/(2n)}\,d\mu
       -n\int\log p\,d\mu .
\end{align*}
All logarithms on the right are integrable: the equilibrium density is
bounded, has square-root decay at its finite positive endpoints, and
\(\omega\) has a positive smooth density there. The prefactor contributes
\(O(n\log q+n\log n)\), and the entropy contributes \(O(n)\).
Division by \(n^2\) therefore yields exactly the value of GEL7 at its
maximizer. This proves GEL10.

\paragraph{GEL5. Finite positive shifts by a proved convex squeeze.}
For each fixed \(n,q\), differentiation under the integral in GEL3 on
every compact subinterval of \(a>-1/2\) gives
\[
 \frac{d}{da}\log J_{n,q}(a)=
    \mathbb E_a\sum_i\log x_i,\qquad
 \frac{d^2}{da^2}\log J_{n,q}(a)=
    \operatorname{Var}_a\!\left(\sum_i\log x_i\right)\geq0.
 \tag{GEL11}
\]
The expectation is with respect to the entire positive integrand in
GEL3 divided by its integral. Extra logarithmic factors are integrable
at zero because \(a>-1/2\), and at infinity by GEL4. Hence the
functions \(F_{n,q}(z)=n^{-2}\log J_{n,q}(nz)\) are differentiable
and convex. Suppose now
\[
 q/n\longrightarrow2,\qquad a/n\longrightarrow2,\qquad
 l>0,\quad l/n\longrightarrow0.
\]
Set
\[
 \mathcal L=\int_{u^2}^{v^2}\log x\,\rho_{2,\pi}(x)\,dx,
 \qquad uK(\kappa)=2,\quad vE(\kappa)=4.
 \tag{GEL12}
\]
For fixed \(0<\eta<2\), both endpoints of the secant from \(a/n\) to
\((a+l)/n\) lie in \((2-\eta/2,2+\eta/2)\) for large \(n\).
Convexity bounds that secant from below by the secant from
\(2-\eta\) to \(2-\eta/2\), and from above by the secant from
\(2+\eta/2\) to \(2+\eta\). Apply GEL10 at these four fixed positive
arguments, then let \(\eta\downarrow0\) and use GEL9. This proves
\[
 \frac{\log J_{n,q}(a+l)-\log J_{n,q}(a)}{nl}
       \longrightarrow\mathcal L.
\]
Combining this with the exact source factors in GEL3 gives
\[
 \boxed{\log\frac{H_n^{(a+l)}}{H_n^{(a)}}
       =2nl\log q+nl\mathcal L+o(nl).}
 \tag{GEL13}
\]
This proof covers moving starting exponents \(a\); it never differentiates
the error in a free-energy asymptotic. The remainder assertion has the
precise scope \(o(nl)\) along the displayed sequences.

\paragraph{GEL6. Both original matrices and the low moment.}
Return to the actual packet
\[
 k=4l+1,\quad e=1+k(m-1),\quad q=e(k+1)^2=2n,
 \qquad m\geq1\text{ fixed},\qquad l\to\infty.
\]
Retain the exact parity identity from HCT35,
\(Z_a=H_n^{(a)}H_{n+1}^{(a)}(H_n^{(a+1)})^2\).
Use GEL13 at the four literal factors, with sizes \(n,n+1,n,n\)
and starting exponents \(q,q,q+1,q+1\), respectively. Every sequence
satisfies its stated domain. Addition gives
\[
 \boxed{\log\frac{Z_{q+l}}{Z_q}
   =2(2q+1)l\log q+(2q+1)l\mathcal L+o(ql).}
 \tag{GEL14}
\]
Thus the contribution from the two original full matrices has been
evaluated, with all four parity blocks still identifiable.

Write \(\mu_a=\int y^{2a}d\sigma(y)\). GEL4 gives, for integer \(a\geq1\),
\[
 \log\mu_a=2a\log a+O(a).
 \tag{GEL15}
\]
One elementary proof of the upper bound uses
\((1+4y^2)^{1/4}\leq1+2y\) for \(y\geq0\); its integrals are
\((2a)!/(\pi/2)^{2a+1}\) and
\(2(2a+1)!/(\pi/2)^{2a+2}\), multiplied by the exact constant
\(2C_\Gamma e^{1/2}\). The bound \(N!\leq N^N\) suffices.
For the lower bound restrict the positive-axis integral to \([a,a+1]\):
there \((1+4y^2)^{-3/4}\geq5^{-3/4}(a+1)^{-3/2}\),
\(y^{2a}\geq a^{2a}\), and
\(e^{-\pi y/2}\geq e^{-\pi(a+1)/2}\). These prove both sides of GEL15
with constants independent of \(a\).

The real-exponent function \(a\mapsto\log\mu_a\) is convex by the
same variance proof as GEL11. Since \(q\) is even and \(1\leq l\leq q\),
its secant on \([q,q+l]\) lies between the secants on \([q/2,q]\)
and \([2q,4q]\). GEL15 bounds both latter secants in absolute value
by \(O(\log q)\). Therefore
\[
 \left|\log\frac{\mu_{q+l}}{\mu_q}\right|=O(l\log q)=o(ql).
 \tag{GEL16}
\]
The original low moment is precisely
\(\mu_{q+l}/\mu_q=m_{2q+2l}/m_{2q}\), so no index is reinterpreted.

\paragraph{GEL7. The leading value of the unchanged signed centre.}
Keep the source product and its exact row intervals:
\[
 P_k=-\sum_{a=q}^{2q-1}\sum_{j=0}^{2l-1}\log(a+j+1/2)
     -\sum_{a=q+1}^{2q}\sum_{j=0}^{2l-1}\log(a+j+1/2).
\]
The two outer sums each have \(q\) terms. Since \(l/q\to0\), direct
Riemann sums, with \(\log\) Lipschitz on \([1,3]\), give
\[
 P_k=-4ql\log q-4ql\int_1^2\log x\,dx+o(ql)
     =-4ql\log q-4ql(2\log2-1)+o(ql).
 \tag{GEL17}
\]
Explicitly, omitting \((j+1/2)/q\) changes each logarithm by at most
\(2l/q\), hence changes the total by at most \(8l^2\); the ordinary
left and right Riemann sums together have error \(O(l)\).

The full finite product calculation WGP10 strengthens this leading
identity to the following enclosure for the very same \(P_k\):
\[
 \begin{split}
 P_k={}&-4lq\log q-4lq(2\log2-1)-4l^2\log2
              +\frac{8l^3+l}{6q}+\mathcal E_{q,l},\\
 -\left\{\frac l{6q}+\frac{8l^3+l}{12q^3}
                 +\frac{8l^4+2l^2}{3q^2}\right\}
 \leq{}&\mathcal E_{q,l}\leq\frac{l^2}{3q^2}.
 \end{split}
 \tag{GEL17a}
\]
Its proof uses the literal midpoint coordinates in WGP3, the signed
quadrature remainders WGP4--5, and the nonnegative cubic remainder
WGP6--9. The subsequent exact refinement WGP12--14 gives
\[
 \begin{split}
 P_k={}&-4lq\log q-4lq(2\log2-1)-4l^2\log2
 +\frac{8l^3+l}{6q}-\frac{l^4+l^2/4}{q^2}
          +\mathcal E^{(3)}_{q,l},\\
 -\frac l{6q}-\frac{8l^3+l}{12q^3}
 \leq{}&\mathcal E^{(3)}_{q,l}\leq
 \frac{l^2}{3q^2}+\frac{8l^4+2l^2}{6q^4}
                 +\frac{7(192l^5+80l^3-2l)}{1440q^3}.
 \end{split}
 \tag{GEL17b}
\]
These original endpoint contributions are retained when the finer
\(q\)-scale calculation is continued. In particular the term
\(-4l^2\log2\) has order \(q\) at \(m=1\), and the cubic term tends
to \(-1/256\) there; the leading limit does not remove either term.

The original quantity is exactly
\(W_k=P_k-\log(\mu_{q+l}/\mu_q)+\log(Z_{q+l}/Z_q)\).
Substitute GEL14, GEL16, and GEL17. The surviving
\(2l\log q\) is \(o(ql)\), so the leading growing-degree calculation is
\[
 \boxed{\frac{W_k}{ql}\longrightarrow
 \mathcal C_\Gamma:=2\mathcal L-4(2\log2-1),\qquad
 \mathcal L=\int_{u^2}^{v^2}\log x\,\rho_{2,\pi}(x)\,dx,\quad
 uK(\kappa)=2,\ vE(\kappa)=4.}
 \tag{GEL18}
\]
Every measure, exponent, matrix dimension, and source factor entering this
limit has been connected to its original object by GEL1--3 and GEL14.
The exact positive density and unique endpoint equations in GEL8 make
the constant an explicitly defined one-dimensional equilibrium integral.
This leading result has remainder \(o(ql)\). It supplies no claim of
an \(o(q)\) remainder, and hence does not erase the next finite-scale
calculation required by the original arithmetic receiver.

The fully proved original-polynomial recurrence estimate RWB20--27
also applies to this identical \(W_k\). It gives the finite bounds
\(lq/64\leq W_k\leq6lq\), and its narrower finite lower remainder
gives \(\liminf W_k/(lq)\geq4\log(27/(8\pi))\). Combining those
proved inequalities with the now proved existence of GEL18 yields
\[
 \boxed{4\log\frac{27}{8\pi}\leq\mathcal C_\Gamma\leq6,
              \qquad\mathcal C_\Gamma>0.}
 \tag{GEL18a}
\]
The strict sign is analytic: RWB22 proves
\(4\log(27/(8\pi))>104/365>0\) using rational integral bounds for
\(\pi\) and the elementary logarithm inequality. Numerical evaluation
of the equilibrium integral is unnecessary for this sign conclusion.

Finally LET20, whose error is \(O(k)=o(q)\) at fixed original
\(m,\delta,\gamma\), transfers the proved limit without changing its sign:
\[
 \frac{\mathfrak T_k}{ql}\longrightarrow\mathcal C_\Gamma,
 \qquad
 \frac{R_k^0-R_k^\sigma}{ql}\longrightarrow-\mathcal C_\Gamma.
 \tag{GEL19}
\]
The arithmetic identity remains
\(B_{\rm ar}=B_0+\delta_\sigma-\mathfrak T_k\), including its original
\(\delta_\sigma\) allowance. Thus GEL19 evaluates the Gamma summand
and its exact signed map into that identity; it makes no unsupported
growth assertion about either of the other two summands.

% END RX INLINE 4


% BEGIN RX INLINE 5; SHA256 9b42c91c8cbc5e8a01aaf130f9d5d5ac3855880bfe85dd72c8ac6add2bfb5479
% Complete new receiving proof; original predecessor files remain preserved.
\section{The evaluated Gamma term in the original arithmetic and cohomology receiver}
\label{sec:WGR}

Retain the actual packet $k=4l+1$, $e=1+k(m-1)$,
$q=e(k+1)^2$, $c=k/2$, $0<\delta<1/2$, $\gamma>2$,
and the original polynomials
\[
